\documentclass[11pt,letterpaper]{article}

% ============================================================
% PAQUETES
% ============================================================
\usepackage[utf8]{inputenc}
\usepackage[T1]{fontenc}
\usepackage{geometry}
\geometry{left=2.7cm,right=2.7cm,top=2.8cm,bottom=2.8cm}
\usepackage{amsmath,amssymb}
\usepackage{graphicx}
\usepackage{booktabs}
\usepackage{longtable}
\usepackage{array}
\usepackage[table]{xcolor}
\usepackage{enumitem}
\usepackage{titlesec}
\usepackage{fancyhdr}
\usepackage[hidelinks,colorlinks=true,linkcolor=black,citecolor=black,urlcolor=blue!50!black]{hyperref}
\usepackage{parskip}
\usepackage{caption}
\usepackage{tcolorbox}
\tcbuselibrary{skins,breakable}

% ============================================================
% COLORES INSTITUCIONALES
% ============================================================
\definecolor{ustablue}{RGB}{0,45,98}
\definecolor{ustagray}{RGB}{90,90,90}
\definecolor{boxbg}{RGB}{237,243,250}

% ============================================================
% FORMATO DE TÍTULOS
% ============================================================
\titleformat{\section}{\normalfont\Large\bfseries\color{ustablue}}{\thesection.}{0.6em}{}
\titleformat{\subsection}{\normalfont\large\bfseries\color{ustablue}}{\thesubsection.}{0.5em}{}
\titleformat{\subsubsection}{\normalfont\normalsize\bfseries\color{ustagray}}{\thesubsubsection.}{0.5em}{}

% ============================================================
% ENCABEZADO / PIE DE PÁGINA
% ============================================================
\pagestyle{fancy}
\fancyhf{}
\renewcommand{\headrulewidth}{0.4pt}
\fancyhead[L]{\small\color{ustagray} Contratación pública y vulnerabilidad socioeconómica municipal}
\fancyhead[R]{\small\color{ustagray} Informe 1 --- Proyecto de grado}
\fancyfoot[C]{\small\thepage}

\setlist[itemize]{leftmargin=1.3em, itemsep=2pt, topsep=3pt}
\setlist[enumerate]{leftmargin=1.5em, itemsep=3pt, topsep=3pt}

\newenvironment{recuadro}[1]{%
  \begin{tcolorbox}[breakable,colback=boxbg,colframe=ustablue,boxrule=0.6pt,arc=1.5pt,
  title={\small\bfseries #1},fonttitle=\small,coltitle=white,colbacktitle=ustablue,
  left=6pt,right=6pt,top=4pt,bottom=4pt]
}{\end{tcolorbox}}

\renewcommand{\contentsname}{Índice}

\begin{document}

% ================================================================
% PORTADA
% ================================================================
\begin{titlepage}
\centering
\vspace*{1.0cm}
{\bfseries\large UNIVERSIDAD SANTO TOMÁS\\
PROYECTO DE GRADO --- ESTADÍSTICA}\\[0.6cm]
{\color{ustablue}\rule{0.85\linewidth}{1pt}}\\[0.5cm]
{\bfseries\Large ESTADO DEL ARTE, MARCO TEÓRICO,\\
DISEÑO METODOLÓGICO Y DIVISIÓN DE TRABAJO\\[0.2cm]
\large INFORME 1 --- VERSIÓN INTEGRADA}\\[0.6cm]
{\color{ustablue}\rule{0.85\linewidth}{1pt}}\\[0.9cm]

{\large\itshape Contratación pública y vulnerabilidad socioeconómica municipal en Colombia}\\[0.9cm]

{\bfseries Pregunta central:}\\[0.15cm]
\parbox{0.82\linewidth}{\centering\itshape
¿Existe una correspondencia entre la distribución territorial de la contratación pública
y las condiciones de vulnerabilidad socioeconómica de los municipios colombianos?}\\[1.3cm]

\begin{tabular}{c}
Juan Diego Gallo Quintero\\
Juan David Parada Fonseca
\end{tabular}\\[1.0cm]

Programa de Estadística --- Octavo semestre\\
Documento de trabajo académico --- Primer informe integrado\\
Septiembre de 2026

\vfill
{\small\itshape Este documento integra, en una sola versión, el estado del arte y
marco teórico del proyecto, el diseño metodológico de muestreo estadístico ya
definido para su ejecución, y la división de trabajo entre los dos integrantes
de la pareja de investigación, articulada con el diagrama de flujo del proyecto.}
\end{titlepage}

\tableofcontents
\newpage

% ================================================================
\section{Alcance y decisiones de investigación}

Este documento constituye la primera versión integrada del estado del arte, el
marco teórico, el diseño metodológico (incluyendo el componente de muestreo
estadístico) y la división de trabajo para un proyecto de grado orientado a
estudiar la correspondencia entre la contratación pública y la vulnerabilidad
socioeconómica de los municipios colombianos. El proyecto toma como punto de
partida el repositorio público \texttt{desarrollo\_social\_y\_economico} de la
Universidad Santo Tomás, pero no lo considera un trabajo propio ni un límite
metodológico: el objetivo es desarrollar una investigación original,
reproducible y estadísticamente rigurosa.

La unidad principal de análisis será el municipio, siempre que la calidad y
disponibilidad de los datos permitan la integración territorial. El nivel
departamental se utilizará como apoyo para contextualización y visualización.
La contratación pública será caracterizada mediante el número de contratos y su
valor monetario, complementados por modalidad de contratación, sector, tipo de
contrato y duración. La vulnerabilidad socioeconómica se conceptualizará a
partir del enfoque multidimensional del DANE y se estudiará la construcción de
un índice sintético basado en el Índice de Pobreza Multidimensional (IPM), con
PCA como técnica estadística central.

El análisis no buscará establecer causalidad. El alcance inferencial se
limitará a estudiar asociación, correspondencia y brechas territoriales. El
análisis espacial será incorporado cuando los diagnósticos de dependencia
espacial lo justifiquen, y el componente de muestreo estadístico --- ya
definido y detallado en la Sección 9 --- se incorpora para resolver dos
problemas metodológicos concretos: la auditoría de calidad de un histórico de
gran volumen y la validación de la estabilidad de las estimaciones.

% ================================================================
\section{Verificación inicial del repositorio y ventana temporal}

La revisión directa del repositorio muestra que su componente de datos
documenta un histórico de procesos SECOP II organizado trimestralmente. El
notebook \texttt{Historico\_de\_Datos.ipynb} registra trimestres desde el
segundo trimestre de 2015 hasta el cuarto trimestre de 2025, aunque este último
trimestre está disponible solo hasta el 5 de octubre en la ejecución
registrada. El histórico reporta 7.827.129 registros y 59 variables, con
última fecha de publicación 2025-10-05 en la ejecución consultada.

Por tanto, la ventana potencial de contratación disponible en el repositorio es
aproximadamente 2015--2025. Sin embargo, esto no significa que todos los años
sean igualmente comparables: será necesario revisar cobertura municipal,
cambios de variables, calidad de identificación territorial y posibles cambios
en el sistema de contratación antes de fijar la ventana definitiva del
estudio --- verificación que, precisamente, se formaliza mediante el muestreo
estratificado de auditoría descrito en la Sección 9.

El README del repositorio indica que la fuente de contratación es SECOP y que
la integración con DANE 2018 fue concebida inicialmente con departamentos y, si
era posible, municipios. También señala que la información de SECOP se
descarga mediante API/CSV y se organiza en archivos Parquet para facilitar el
tratamiento de grandes volúmenes.

\begin{recuadro}{Decisión metodológica}
Se utilizará 2015--2025 como ventana de exploración del componente contractual,
pero el periodo analítico definitivo se seleccionará después de auditar la
cobertura y armonización de las variables socioeconómicas. Esta decisión evita
forzar una comparabilidad temporal que los datos de DANE podrían no sostener.
\end{recuadro}

% ================================================================
\section{Estado del arte}

\subsection{Desigualdad territorial, gasto público y desarrollo regional}

La literatura colombiana ha mostrado históricamente que las diferencias
territoriales en ingreso, capacidad fiscal y prestación de servicios son
persistentes. Ardila-Rueda (2004) estudió la convergencia del ingreso per
cápita departamental y el papel del consumo e inversión pública, encontrando
una alta persistencia en la distribución del ingreso y efectos del gasto
público --- especialmente de la inversión --- sobre la posición relativa de
algunos departamentos, sin que ello implicara una transformación completa de
la distribución territorial.

Bonet-Morón (2006) analizó los desequilibrios de la descentralización y mostró
diferencias importantes en ingresos fiscales per cápita y en las condiciones de
provisión de servicios entre entidades territoriales. Su argumento es
especialmente relevante para este proyecto porque advierte que la
descentralización y las transferencias no garantizan por sí mismas equidad
territorial.

Estos antecedentes establecen una primera base conceptual: la distribución de
recursos públicos debe analizarse en relación con las características
estructurales de los territorios y no solamente en términos de montos
absolutos. Para el presente proyecto, esto justifica la utilización de medidas
per cápita y medidas relativas, además del valor total de contratación.

\subsection{Pobreza multidimensional y vulnerabilidad territorial}

El enfoque de pobreza multidimensional constituye una referencia central para
operacionalizar la vulnerabilidad socioeconómica. El DANE adopta una medición
multidimensional que integra dimensiones y privaciones relacionadas con las
condiciones de vida. A nivel municipal, la fuente censal de 2018 permite
disponer de una medida territorial de IPM con amplia cobertura municipal.

El uso del IPM resulta especialmente pertinente porque permite evitar una
definición unidimensional de vulnerabilidad basada exclusivamente en ingreso.
No obstante, el proyecto va un paso más allá: el IPM será el núcleo conceptual
del indicador, mientras que variables adicionales disponibles y compatibles se
evaluarán mediante PCA para construir un índice sintético de vulnerabilidad
socioeconómica municipal.

Un antecedente reciente particularmente cercano al componente territorial es
Monsalvo Herrera y Jiménez (2025), quienes analizaron el IPM en municipios de
Cundinamarca mediante indicadores locales de asociación espacial (LISA). El
estudio encontró agrupamientos espaciales de pobreza y plantea la relevancia
del efecto vecindario para orientar acciones de política pública. Este
antecedente es directamente útil para justificar que, si los diagnósticos lo
respaldan, el proyecto incorpore estadística espacial.

\subsection{Contratación pública como instrumento de política territorial}

La contratación pública puede interpretarse no solo como una actividad
administrativa de compra de bienes y servicios, sino como un mecanismo mediante
el cual el Estado canaliza recursos hacia territorios, sectores y proveedores.
La literatura internacional reciente ha estudiado explícitamente esta dimensión
territorial. Ferry et al. (2023) analizaron más de 90.000 contratos de
autoridades locales del Reino Unido y encontraron diferencias regionales en la
selección de proveedores, planteando que la contratación puede tener una
dimensión territorial asociada al desarrollo regional.

Arachia et al. (2024), a partir de más de 50.000 contratos de obras públicas
municipales en Italia, estudiaron la cooperación intermunicipal mediante
modelos de efectos fijos y estimadores de \textit{matching}. Sus resultados
muestran que la forma institucional de organizar la contratación puede
relacionarse con resultados de ejecución. Aunque el contexto institucional es
distinto al colombiano, el estudio es útil para mostrar que la contratación
municipal puede analizarse estadísticamente como fenómeno territorial y no
únicamente jurídico.

Sarter (2024) plantea una relación multidimensional entre contratación pública
y desigualdad laboral, argumentando que el objeto comprado, los requisitos del
proceso y el tipo de proveedor pueden afectar resultados distributivos. Este
marco apoya la decisión de no limitar el análisis a los montos, sino incorporar
modalidad, sector, tipo de contrato y características de ejecución.

Rodriguez-Plesa et al. (2022) muestran, desde la contratación pública
sostenible a nivel local, que las compras gubernamentales pueden utilizarse
estratégicamente para abordar objetivos sociales, económicos y ambientales.
Para el presente proyecto, esto fortalece la interpretación de la contratación
como instrumento potencial de política pública, aunque no permite asumir que
todo gasto contractual tenga un efecto social positivo.

\subsection{SECOP II y análisis estadístico de datos abiertos en Colombia}

La disponibilidad de SECOP II permite estudiar la contratación pública con
microdatos administrativos. El conjunto de contratos electrónicos identifica
información sobre entidades, ubicación, fechas, sectores, valores y otras
características del contrato. El repositorio estudiado utiliza esta fuente y ha
construido un proceso de extracción trimestral y consolidación histórica.

En Colombia, Figueroa-Gómez y Galpin (2024) utilizaron datos abiertos de SECOP
II para desarrollar modelos de clasificación destinados a identificar procesos
de contratación de interés, señalando que la investigación aplicada a
contratación pública colombiana todavía se encuentra en una etapa de
desarrollo. Salazar, Pérez y Gallego (2024) desarrollaron VigIA, una
herramienta basada en datos abiertos de SECOP II para detectar ineficiencias y
construir índices de riesgo; su unidad de análisis y objetivo se concentran en
riesgo e ineficiencia contractual, no en correspondencia socioeconómica
territorial. Gutiérrez Vanegas (2024) utilizó \textit{Isolation Forest} sobre
datos abiertos de SECOP II para identificar banderas rojas asociadas con
periodos inusualmente cortos, valores anómalos y procesos con una sola oferta;
nuevamente, el problema de investigación es la detección de anomalías y no la
equidad territorial.

Díaz Díez (2023) analiza el carácter transaccional de SECOP II y su relación
con los principios de publicidad y transparencia. Desde la perspectiva del
proyecto, esta literatura ayuda a comprender la naturaleza institucional de la
fuente y sus limitaciones, evitando interpretar los datos de SECOP como una
medición directa del desarrollo social.

\subsection{Estadística espacial y brechas territoriales}

La dimensión espacial es relevante porque municipios próximos pueden compartir
condiciones económicas, sociales, geográficas e institucionales. El estudio de
Monsalvo Herrera y Jiménez (2025) evidencia que el IPM municipal puede
presentar patrones de dependencia espacial y agrupamientos identificables
mediante LISA. Esto sugiere que la correlación convencional entre
vulnerabilidad y contratación podría ser insuficiente si existe
autocorrelación espacial.

En consecuencia, el análisis espacial no se asumirá automáticamente como una
técnica decorativa. La secuencia metodológica será: primero, diagnóstico de
autocorrelación espacial; después, si existe evidencia suficiente,
incorporación de modelos o indicadores espaciales. Esta regla es coherente con
una práctica estadística rigurosa y evita sobreajustar el proyecto a técnicas
avanzadas sin necesidad.

\subsection{Síntesis crítica y vacío de investigación}

La revisión muestra cuatro cuerpos de literatura parcialmente conectados.
Primero, la investigación sobre desigualdad regional colombiana documenta
disparidades persistentes y diferencias en capacidad fiscal y prestación de
servicios. Segundo, la literatura sobre pobreza multidimensional permite
operacionalizar condiciones de vulnerabilidad con una perspectiva territorial.
Tercero, la literatura sobre contratación pública reconoce que la compra
estatal puede ser un instrumento de política y que sus características pueden
tener implicaciones distributivas. Cuarto, la literatura reciente sobre SECOP
II demuestra el potencial de los datos abiertos para estadística, aprendizaje
automático y control de riesgos.

Sin embargo, en la revisión inicial no se identificó un estudio que integre de
manera sistemática, a escala nacional y con unidad municipal, un índice
multidimensional de vulnerabilidad construido desde el IPM/PCA con medidas de
contratación pública y un análisis conjunto de asociación, correspondencia y
brechas territoriales. Los trabajos encontrados sobre SECOP II se concentran
principalmente en riesgo, eficiencia, clasificación o detección de anomalías;
los trabajos sobre pobreza territorial estudian la distribución y dependencia
espacial de la pobreza, pero no la correspondencia con la distribución de
contratos públicos.

Este vacío es provisional y deberá comprobarse mediante una revisión
bibliográfica más exhaustiva en Scopus, Web of Science, SciELO, Redalyc y
repositorios universitarios. No obstante, constituye una justificación
razonable para la pregunta central del proyecto.

% ================================================================
\section{Matriz bibliográfica inicial}

\begin{center}
\footnotesize
\begin{longtable}{p{2.1cm}p{3.0cm}p{3.0cm}p{3.6cm}p{3.6cm}}
\toprule
\textbf{Referencia} & \textbf{Datos y unidad} & \textbf{Método} & \textbf{Resultado principal} & \textbf{Limitación / utilidad} \\
\midrule
\endfirsthead
\toprule
\textbf{Referencia} & \textbf{Datos y unidad} & \textbf{Método} & \textbf{Resultado principal} & \textbf{Limitación / utilidad} \\
\midrule
\endhead
Ardila-Rueda (2004) & Ingreso departamental + gasto público; Colombia & Análisis de dinámica distributiva & Persistencia regional; la inversión afecta posiciones relativas & Antiguo, nivel departamental; fundamenta relación gasto--territorio \\
Bonet-Morón (2006) & Ingresos fiscales, transferencias, servicios; departamentos/municipios & Análisis fiscal regional & Persisten desigualdades fiscales y de servicios & Antiguo; justifica medidas relativas y equidad territorial \\
Rodriguez-Plesa et al. (2022) & Contratación sostenible local; EE.\,UU./internacional & Indicadores/modelos institucionales & La compra pública puede perseguir objetivos sociales/económicos & No estudia vulnerabilidad municipal; marco de contratación como política pública \\
Ferry et al. (2023) & 90.000 contratos 2015--2019; autoridades locales, Reino Unido & Análisis cuantitativo regional & Existen patrones territoriales de selección de proveedores & Contexto institucional distinto; aporta dimensión territorial de contratación \\
Sarter (2024) & Literatura sobre contratación; procesos/proveedores & Marco conceptual & Objeto, proceso y proveedor pueden afectar desigualdad laboral & No es estudio municipal colombiano; justifica variables contractuales \\
Arachia et al. (2024) & 50.905 contratos de obras; municipios, Italia & Efectos fijos + \textit{matching} & Cooperación intermunicipal se relaciona con desempeño & No estudia pobreza; referencia para panel/municipios \\
Figueroa-Gómez \& Galpin (2024) & SECOP II; procesos de contratación, Colombia & Modelos de clasificación & Potencial de ML sobre contratación pública & Enfoque en licitaciones de interés; antecedente de ciencia de datos \\
Salazar et al. (2024) & SECOP II + sanciones; contratos, Colombia & ML + índices de riesgo & Detección de ineficiencias y riesgos & No estudia vulnerabilidad territorial; referencia metodológica para índices \\
Gutiérrez Vanegas (2024) & SECOP II; contratos, Colombia & \textit{Isolation Forest} & Detección de banderas rojas & Trabajo de grado, foco en anomalías; referencia de ML aplicado a SECOP II \\
Díaz Díez (2023) & SECOP II y normativa; sistema de contratación, Colombia & Análisis jurídico & SECOP II es transaccional y fortalece publicidad/transparencia & No es análisis estadístico; marco institucional de la fuente \\
Monsalvo Herrera \& Jiménez (2025) & IPM municipal; municipios de Cundinamarca, Colombia & LISA/análisis espacial & Agrupamientos espaciales de pobreza & Cobertura regional; justifica diagnóstico espacial \\
DANE (2018/2024) & CNPV/IPM; municipios, Colombia & Medición multidimensional & Base oficial de condiciones de vida & Cobertura temporal no anual; núcleo del índice de vulnerabilidad \\
\bottomrule
\end{longtable}
\end{center}

% ================================================================
\section{Marco teórico}

\subsection{Desarrollo territorial y desigualdad}

El desarrollo territorial se entiende aquí como un proceso diferenciado
espacialmente en el que interactúan condiciones económicas, sociales,
institucionales y de provisión de bienes públicos. La desigualdad territorial
implica que los municipios enfrentan capacidades y necesidades diferentes, por
lo que la comparación de recursos públicos debe considerar el tamaño
poblacional y las condiciones socioeconómicas.

\subsection{Pobreza multidimensional y vulnerabilidad socioeconómica}

El IPM proporciona una base conceptual para reconocer que el bienestar no puede
reducirse al ingreso. Para este proyecto, vulnerabilidad socioeconómica será un
concepto operativo más amplio que pobreza: representará un conjunto de
privaciones o condiciones estructurales que pueden limitar el desarrollo de la
población municipal. El IPM será el referente principal y las variables
adicionales solo se incorporarán cuando exista respaldo conceptual, calidad
estadística y disponibilidad comparable.

\subsection{Índices sintéticos y PCA}

Un índice sintético busca resumir múltiples variables correlacionadas en una
medida interpretable. El PCA permite transformar un conjunto de variables
estandarizadas en componentes ortogonales que maximizan la varianza explicada.
Si $Z$ representa la matriz estandarizada de indicadores, la descomposición
puede expresarse como:
\[
Z = U \Lambda V'
\]

La construcción del índice no se decidirá únicamente por la mayor varianza
explicada. Se evaluarán correlaciones, adecuación muestral, cargas, signos,
estabilidad y coherencia sustantiva. También se invertirán las variables cuyo
aumento represente menor vulnerabilidad, para asegurar que la dirección del
índice sea consistente. La estabilidad del índice --- signos, magnitudes
relativas y varianza explicada --- se valida formalmente mediante remuestreo
\textit{bootstrap}, descrito en la Sección 9.4.

\subsection{Contratación pública como distribución de recursos}

Para el análisis, la contratación pública se operacionalizará mediante al
menos cuatro familias de indicadores: magnitud monetaria, frecuencia
contractual, estructura contractual y temporalidad. Entre los indicadores
candidatos se encuentran valor total, número de contratos, valor per cápita,
valor mediano, ticket promedio, participación sectorial, distribución por
modalidad y duración.

El valor per cápita será especialmente importante porque un valor absoluto
elevado puede reflejar simplemente mayor población. La comparación entre
municipios requiere, por tanto, medidas que controlen al menos parcialmente el
tamaño poblacional.

\subsection{Asociación, correspondencia y brechas}

La asociación se refiere a la relación estadística entre vulnerabilidad y
contratación. Se utilizarán correlaciones de Pearson o Spearman según escala,
distribución y supuestos, complementadas con modelos multivariados y, como
verificación no paramétrica, con la prueba de permutación descrita en la
Sección 9.5.

La correspondencia será un concepto analítico distinto: no preguntará
solamente si dos variables se mueven juntas, sino si la intensidad de
contratación resulta proporcional o compatible con la posición relativa de
vulnerabilidad de los municipios. Se estudiarán diferentes definiciones de
correspondencia y se seleccionará aquella que pueda justificarse
estadísticamente.

Las brechas territoriales identificarán municipios cuyo nivel de contratación
observado sea bajo o alto respecto de lo esperado según su vulnerabilidad y
otras características relevantes. El lenguaje será deliberadamente prudente:
una brecha no significará automáticamente ineficiencia, corrupción o mala
política pública.

\subsection{Dependencia espacial}

La estadística espacial será utilizada si existe evidencia de dependencia
espacial. Se evaluará Moran global y, posteriormente, indicadores locales
LISA. Si la dependencia espacial afecta los residuos de un modelo, podrán
considerarse modelos espaciales como SAR/lag, SEM o SDM, según diagnóstico y
factibilidad.

\subsection{Estructura temporal y datos de panel}

La disponibilidad de contratación desde 2015 hasta 2025 permite explorar una
estructura municipio--año. Cuando las variables socioeconómicas no tengan
frecuencia anual, no se imputará artificialmente su evolución; la solución se
definirá después de auditar la disponibilidad de indicadores. Si existe
suficiente variación temporal, los modelos de panel con efectos fijos o
aleatorios se evaluarán mediante pruebas y criterios de especificación.

% ================================================================
\section{Modelo conceptual de investigación}

La estructura conceptual del proyecto es:

\begin{center}
\itshape
Condiciones socioeconómicas $\rightarrow$ Índice de Vulnerabilidad
$\rightarrow$ Indicadores de Contratación $\rightarrow$ Asociación
$\rightarrow$ Correspondencia $\rightarrow$ Brechas $\rightarrow$
Diagnóstico espacial/temporal $\rightarrow$ Clasificación territorial.
\end{center}

La pregunta no presupone que los municipios más vulnerables deban recibir
necesariamente más contratos. La investigación pretende determinar
empíricamente si existe un patrón de correspondencia y cómo se distribuyen las
desviaciones respecto de ese patrón. El muestreo estratificado de auditoría y
el \textit{bootstrap} de validación (Sección 9) se insertan en los dos primeros
eslabones de esta cadena --- calidad del dato de entrada y estabilidad del
índice --- antes de avanzar a las etapas de asociación y correspondencia.

% ================================================================
\section{Variables y operacionalización preliminar}

\begin{center}
\footnotesize
\begin{longtable}{p{2.2cm}p{2.6cm}p{4.6cm}p{1.8cm}p{2.6cm}}
\toprule
\textbf{Constructo} & \textbf{Variable} & \textbf{Operacionalización} & \textbf{Escala} & \textbf{Rol} \\
\midrule
\endfirsthead
\toprule
\textbf{Constructo} & \textbf{Variable} & \textbf{Operacionalización} & \textbf{Escala} & \textbf{Rol} \\
\midrule
\endhead
Vulnerabilidad & Índice sintético & PCA basado en IPM + indicadores compatibles & Continua & Principal \\
IPM & Pobreza multidimensional & Fuente DANE & Continua & Base conceptual \\
Contratación & Valor contratado & Suma por municipio--año & Continua & Dependiente/resultado \\
Contratación & Número de contratos & Conteo por municipio--año & Conteo & Resultado \\
Contratación & Valor per cápita & Valor contratado / población & Continua & Resultado \\
Estructura & Modalidad & Frecuencia/participación por modalidad & Categórica & Explicativa/descriptiva \\
Estructura & Sector & Participación por sector & Categórica & Explicativa/descriptiva \\
Estructura & Tipo de contrato & Participación por tipo & Categórica & Explicativa/descriptiva \\
Temporal & Duración & Días entre inicio y terminación, cuando exista & Continua & Explicativa/descriptiva \\
Contexto & Población & Población municipal & Continua & Control \\
Territorio & Municipio/departamento & DIVIPOLA u homologación oficial & Categórica & Identificación \\
Calidad del dato & Indicadores de auditoría & Duplicidad, inconsistencia territorial, campos incompletos, atipicidad (ver Sección 9.3) & Binaria/proporción & Control de calidad \\
\bottomrule
\end{longtable}
\end{center}

% ================================================================
\section{Hipótesis de trabajo}

\begin{recuadro}{Hipótesis H1 -- H6}
\textbf{H1.} Existe una asociación estadísticamente significativa entre el
índice de vulnerabilidad socioeconómica municipal y la intensidad de
contratación pública.\\[4pt]
\textbf{H2.} La asociación entre vulnerabilidad y contratación difiere según la
forma de medir la contratación: valor absoluto, valor per cápita y número de
contratos.\\[4pt]
\textbf{H3.} Existen municipios cuya intensidad de contratación presenta
desviaciones sistemáticas respecto de la correspondencia esperada con su
vulnerabilidad socioeconómica.\\[4pt]
\textbf{H4.} La distribución espacial de vulnerabilidad y/o contratación
presenta autocorrelación espacial cuando los diagnósticos correspondientes así
lo indiquen.\\[4pt]
\textbf{H5.} La tasa de inconsistencias (duplicados, campos incompletos,
valores atípicos no justificados) en el histórico de SECOP II, estimada
mediante muestreo estratificado, no supera un umbral que comprometa la validez
de las métricas de contratación agregadas a nivel municipio--año.\\[4pt]
\textbf{H6.} Las cargas factoriales del componente principal utilizado para
construir el IVS son estables ante remuestreo \textit{bootstrap}: sus
intervalos de confianza no cambian de signo ni de orden de magnitud relativo
entre indicadores.
\end{recuadro}

Ninguna de estas hipótesis plantea causalidad. Los resultados se interpretarán
como asociación, correspondencia y validez estadística.

% ================================================================
\section{Diseño metodológico: integración de muestreo estadístico}

Esta sección formaliza el componente de muestreo estadístico del proyecto. No
se trata de una posibilidad exploratoria: es el diseño ya definido que se
ejecutará como parte del pipeline metodológico, en los puntos precisos del
diagrama de flujo que se indican en cada procedimiento.

\subsection{Naturaleza poblacional de los datos y necesidad del muestreo}

Tanto SECOP II (7.827.129 registros = universo administrativo de contratos, no
una muestra de ellos) como el IPM del DANE 2018 (medición de fuente censal) son
datos de naturaleza \textbf{poblacional}. El muestreo no se introduce para
reemplazar el análisis censal, sino para resolver dos problemas concretos que
la propia escala de los datos genera:

\begin{enumerate}
\item La auditoría manual exhaustiva de 7,8 millones de registros para
diagnosticar su calidad es inviable en el tiempo disponible para el proyecto.
\item La construcción del IVS mediante PCA y la medida de correspondencia
producen estimaciones puntuales cuya estabilidad no queda establecida solo con
los supuestos asintóticos convencionales.
\end{enumerate}

\begin{recuadro}{Principio metodológico}
El muestreo se utiliza para (a) hacer viable la auditoría de calidad de una
base de gran volumen mediante inferencia hacia esa misma población, y (b)
cuantificar la incertidumbre de estimaciones (PCA, correspondencia) mediante
remuestreo, cuando la teoría asintótica clásica es insuficiente o sus
supuestos son dudosos.
\end{recuadro}

\subsection{Justificación del diseño muestral seleccionado: por qué muestreo aleatorio estratificado}

La condición determinante para elegir el diseño muestral es que la calidad de
los registros de SECOP II \textbf{no se distribuye de manera homogénea} a lo
largo de la ventana 2015--2025. Existen al menos tres fuentes documentadas de
heterogeneidad conocida \textit{a priori}: (i) cambios normativos en las reglas
de reporte de la contratación pública electrónica a lo largo del periodo; (ii)
migraciones y actualizaciones de la plataforma SECOP II, que históricamente han
generado periodos de transición con mayor incidencia de campos vacíos,
duplicados o cambios de esquema; y (iii) crecimiento no lineal del volumen de
contratos publicados año a año.

Esta heterogeneidad conocida y estructurada entre años es, precisamente, la
condición bajo la cual la teoría de muestreo establece que la estratificación
domina a las demás alternativas en eficiencia estadística (Cochran, 1977;
Lohr, 2019). La siguiente tabla sustenta, frente a cada alternativa clásica,
por qué el diseño estratificado es la opción seleccionada y no otra:

\begin{center}
\small
\begin{longtable}{p{2.5cm}p{5.9cm}p{4.8cm}}
\toprule
\textbf{Diseño} & \textbf{Por qué no se selecciona} & \textbf{Criterio técnico} \\
\midrule
\endfirsthead
\toprule
\textbf{Diseño} & \textbf{Por qué no se selecciona} & \textbf{Criterio técnico} \\
\midrule
\endhead
Aleatorio simple (MAS) & No garantiza representación proporcional de los años/periodos con mayor incidencia de error; puede sobre-representar por azar los años de mayor volumen y dejar sin observaciones los años de transición de plataforma, que son los más críticos para el diagnóstico. & Varianza del estimador MAS $\geq$ varianza del estimador estratificado siempre que exista heterogeneidad entre estratos; aquí la heterogeneidad está documentada. \\
Sistemático & Requiere un marco muestral ordenado y un salto fijo $k$. Si el archivo está ordenado por fecha o identificador correlativo, el salto puede coincidir con patrones periódicos de carga de datos, introduciendo sesgo de periodicidad. & Riesgo de covarianza no nula entre el intervalo $k$ y ciclos administrativos de reporte. \\
Conglomerados (p.\ ej.\ por entidad contratante) & Auditar conglomerados completos implica revisar registros muy similares entre sí (una misma entidad repite sus propios errores de captura), lo cual reduce la información efectiva por unidad auditada. & Efecto de diseño ($\textit{design effect}$) $>1$ cuando existe correlación intraconglomerado. \\
\textbf{Estratificado (seleccionado)} & Aprovecha la heterogeneidad conocida (año/periodo de plataforma) para reducir la varianza, garantiza representación mínima de cada año --- incluidos los años críticos de transición --- y entrega una estimación de calidad por estrato. & Varianza del estimador estratificado $\leq$ varianza MAS bajo asignación proporcional, y estrictamente menor bajo asignación de Neyman (Cochran, 1977, cap.\ 5). \\
\bottomrule
\end{longtable}
\end{center}

Adicionalmente, únicamente un diseño estratificado produce, sin supuestos
adicionales ni post-estratificación correctiva, estimaciones separadas y
válidas \textbf{por año}, que se usarán directamente para priorizar las reglas
de limpieza automatizada por periodo en la Fase 3. Por estas razones ---
heterogeneidad documentada, menor varianza y utilidad directa para la limpieza
posterior --- se selecciona el muestreo aleatorio estratificado como el diseño
del proyecto.

\subsection{Procedimiento 1: auditoría de calidad de SECOP II mediante muestreo estratificado}

\textbf{Objetivo.} Estimar, con un nivel de confianza y margen de error
predefinidos, la proporción de registros de SECOP II que presentan problemas
de calidad, sin necesidad de revisar el universo completo de 7,8 millones de
contratos.

\textbf{Marco muestral y unidad de muestreo.} Población objetivo: totalidad de
registros de SECOP II del histórico consolidado, dentro de la ventana temporal
definitiva que se fije tras esta misma auditoría (provisionalmente
2015--2025). Unidad de muestreo: el contrato individual. Marco muestral: el
archivo histórico en formato Parquet ya consolidado por trimestre, una vez
verificada su completitud a nivel de archivo.

\textbf{Pasos de ejecución.}
\begin{enumerate}
\item Definir los estratos $h = 1, \dots, L$ cruzando año (o trimestre) con
rango de valor del contrato (cuartiles).
\item Calcular $N_h$, el número de registros en cada estrato dentro del marco
muestral.
\item Determinar el tamaño de muestra total $n$ a partir de una proporción
esperada de inconsistencia $p$ (usar $p=0.5$ como escenario conservador si no
hay estimación previa), un margen de error $e$ y un nivel de confianza
$1-\alpha$:
\[
n_0 = \frac{z_{1-\alpha/2}^2\, p(1-p)}{e^2}, \qquad
n = \frac{n_0}{1 + \dfrac{n_0 - 1}{N}}
\]
aplicando la corrección por población finita, relevante porque $N$ es grande
pero conocido y fijo.
\item Asignar la muestra entre estratos mediante asignación proporcional
como opción base ($n_h = n\cdot N_h/N$), evaluando asignación de Neyman como
alternativa en estratos con mayor variabilidad esperada de calidad:
\[
n_h^{\text{Neyman}} = n \cdot \frac{N_h S_h}{\sum_h N_h S_h}
\]
\item Extraer la muestra mediante muestreo aleatorio simple dentro de cada
estrato, con semilla aleatoria documentada para reproducibilidad.
\item Auditar cada registro muestreado sobre cinco dimensiones binarias de
calidad: duplicidad; inconsistencia territorial (código DIVIPOLA inexistente o
contradictorio); campos clave incompletos; valor atípico no justificado; e
incoherencia temporal (duración negativa o implausible).
\item Estimar, para cada dimensión, la proporción poblacional y su intervalo de
confianza mediante el estimador estratificado:
\[
\hat p_{st} = \sum_{h=1}^{L} \frac{N_h}{N}\,\hat p_h, \qquad
\widehat{\text{Var}}(\hat p_{st}) = \sum_{h=1}^{L}
\left(\frac{N_h}{N}\right)^2 \frac{\hat p_h (1-\hat p_h)}{n_h - 1}
\left(1-\frac{n_h}{N_h}\right)
\]
\item Redactar el reporte de calidad y la recomendación de reglas de limpieza
priorizadas por estrato, según la tasa de error estimada.
\end{enumerate}

\textbf{Entregables.} Ficha técnica del diseño muestral; base auditada con las
cinco dimensiones de calidad codificadas; reporte de estimaciones poblacionales
con intervalos de confianza por dimensión y por estrato; recomendación de
reglas de limpieza para la Fase 3.

\textbf{Ubicación en el flujo del proyecto.} Este procedimiento se ejecuta
entre las fases 2 (Revisión de datos y fuentes) y 3 (Preparación y limpieza de
datos); formaliza estadísticamente el paso 2.1 (auditoría de la base SECOP II)
ya previsto.

\subsection{Procedimiento 2: validación del Índice de Vulnerabilidad mediante \textit{bootstrap}}

\textbf{Objetivo.} Cuantificar la estabilidad estadística de los resultados
del PCA (Fase 4) utilizado para construir el IVS, complementando la validación
por adecuación muestral (KMO/Bartlett) con una medida explícita de
incertidumbre sobre las cargas factoriales.

\textbf{Justificación: por qué \textit{bootstrap} y no otra alternativa.} El
PCA entrega cargas factoriales y varianza explicada como estimaciones
puntuales; sin una medida de incertidumbre no es posible saber si esas cargas
son estables o son un artefacto de la muestra particular de municipios
disponible. Existen al menos tres alternativas para cuantificar esa
incertidumbre, y se descartan por razones concretas:

\begin{center}
\small
\begin{longtable}{p{2.8cm}p{5.9cm}p{4.8cm}}
\toprule
\textbf{Alternativa} & \textbf{Por qué no se selecciona} & \textbf{Criterio técnico} \\
\midrule
\endfirsthead
\toprule
\textbf{Alternativa} & \textbf{Por qué no se selecciona} & \textbf{Criterio técnico} \\
\midrule
\endhead
Fórmulas analíticas asintóticas (teoría clásica de PCA) & Suponen normalidad multivariada de los indicadores y un tamaño de muestra grande para que la aproximación asintótica sea válida. Los indicadores de pobreza y vulnerabilidad son característicamente asimétricos (colas pesadas, concentración en valores bajos), por lo que el supuesto de normalidad es cuestionable con $\sim$1.100 municipios. & La varianza asintótica de los \textit{eigenvectors} depende de la separación entre \textit{eigenvalues}; con valores propios cercanos, la aproximación de primer orden es inestable, mientras que el \textit{bootstrap} no depende de esa aproximación. \\
\textit{Jackknife} (\textit{leave-one-out}) & Es conocido que el \textit{jackknife} se comporta erráticamente para estadísticos no suaves como los vectores propios de un PCA, especialmente cuando dos \textit{eigenvalues} están próximos (riesgo de rotación/cambio de signo abrupto al remover una sola observación). & El \textit{jackknife} solo genera $N$ pseudo-réplicas (una por municipio removido) frente a $B\geq 1000$ del \textit{bootstrap}, ofreciendo una aproximación mucho más pobre de la distribución muestral completa. \\
Validación cruzada (\textit{cross-validation}) & Está diseñada para evaluar capacidad predictiva o error de generalización, no para construir intervalos de confianza de parámetros como cargas factoriales o varianza explicada. Usarla aquí respondería una pregunta distinta a la que exige H6. & No genera una distribución muestral del parámetro de interés, solo una estimación de error de predicción fuera de muestra. \\
\textbf{\textit{Bootstrap} (seleccionado)} & No requiere supuestos distribucionales sobre los indicadores: remuestrea directamente la distribución empírica observada. Genera $B\geq 1000$ réplicas completas del PCA, permitiendo construir intervalos de confianza (percentil o BC$_a$) que capturan asimetría real en la distribución de las cargas, en lugar de asumir simetría como la teoría asintótica. & Es el método estándar en estadística multivariada para validar la estabilidad de cargas factoriales e índices sintéticos basados en PCA (Efron \& Tibshirani, 1993). \\
\bottomrule
\end{longtable}
\end{center}

\textbf{Pasos de ejecución.}
\begin{enumerate}
\item Definir la matriz de indicadores estandarizados $Z$ (municipios $\times$
variables IPM/compatibles), ya construida en la Fase 4.
\item Generar $B \geq 1000$ remuestras \textit{bootstrap} por remuestreo con
reemplazo de las filas (municipios) de $Z$; cuando exista agrupamiento
departamental relevante, ejecutar adicionalmente un \textit{bootstrap} por
bloques (remuestreo de departamentos completos) para no romper la posible
dependencia espacial intradepartamental.
\item Para cada remuestra $b=1,\dots,B$, reajustar el PCA y registrar: las
cargas factoriales del/los componente(s) retenido(s), la proporción de
varianza explicada y el IVS resultante por municipio.
\item Construir, para cada carga factorial y para la varianza explicada, un
intervalo de confianza \textit{bootstrap} (percentil o BC$_a$).
\item Verificar H6: comprobar si el signo y el orden relativo de magnitud de
cada carga se mantienen estables a través de las $B$ remuestras.
\end{enumerate}

\textbf{Criterio de interpretación.} Cargas con intervalos que no cruzan cero y
mantienen signo consistente: variable con contribución estable al índice.
Cargas con intervalos que cruzan cero o cambian de signo: variable candidata a
revisión. Municipios cuyo IVS varía sustancialmente entre remuestras: señal de
sensibilidad del índice para ese municipio.

\textbf{Entregables.} Intervalos de confianza \textit{bootstrap} para cargas
factoriales y varianza explicada; gráfico de distribución \textit{bootstrap}
del IVS para los municipios en los extremos del ranking; reporte de
sensibilidad del índice.

\textbf{Ubicación en el flujo del proyecto.} Se ejecuta dentro de la Fase 4.6
(Validación y sensibilidad del índice), ya prevista en el flujo original; este
procedimiento especifica el método concreto para ejecutar esa validación.

\subsection{Procedimiento 3: prueba de permutación para la correspondencia vulnerabilidad--contratación}

\textbf{Objetivo.} Aportar una vía de inferencia no paramétrica,
complementaria a las pruebas clásicas de correlación y a los modelos de
regresión de la Fase 7, para evaluar si la asociación observada entre el IVS y
los indicadores de contratación es mayor de lo esperable bajo aleatoriedad.

\textbf{Justificación: por qué prueba de permutación y no otra alternativa.}
La pregunta que exige H1 es puntual: ¿la asociación observada entre IVS y
contratación es mayor de lo esperable bajo aleatoriedad? Existen varias formas
de responderla, y se descartan por razones concretas frente a la permutación:

\begin{center}
\small
\begin{longtable}{p{2.8cm}p{5.9cm}p{4.8cm}}
\toprule
\textbf{Alternativa} & \textbf{Por qué no se selecciona} & \textbf{Criterio técnico} \\
\midrule
\endfirsthead
\toprule
\textbf{Alternativa} & \textbf{Por qué no se selecciona} & \textbf{Criterio técnico} \\
\midrule
\endhead
Prueba paramétrica (correlación de Pearson / regresión OLS) & Supone normalidad bivariada y relación lineal entre IVS y contratación. Los indicadores de contratación (valor total, valor per cápita) son característicamente asimétricos a la derecha, con una masa importante de municipios en valores bajos o nulos y una cola larga de valores altos, lo que viola el supuesto de normalidad requerido para que el valor $p$ paramétrico sea válido. & El valor $p$ paramétrico depende de la distribución teórica t de Student bajo normalidad; con datos sesgados, esa distribución teórica no coincide con la distribución real del estadístico, sesgando la inferencia. \\
Correlación de Spearman/Kendall con $p$ asintótico & El valor $p$ reportado por defecto en la mayoría de \textit{software} es una aproximación asintótica (válida solo para muestras grandes) y pierde precisión cuando hay muchos empates (\textit{ties}) --- frecuentes aquí porque varios municipios comparten el mismo número de contratos o el mismo rango de IVS. & La aproximación asintótica de Spearman/Kendall no está diseñada para corregir empates de forma exacta; la permutación sí, porque recalcula la distribución nula empírica directamente sobre los datos observados, empates incluidos. \\
Prueba de bondad de ajuste (chi-cuadrado) sobre categorías de IVS/contratación & Exige discretizar variables continuas en categorías, lo que descarta información al convertir el IVS y la contratación en pocos grupos, y el resultado depende arbitrariamente de dónde se coloquen los puntos de corte. & Pérdida de potencia estadística frente a un estadístico calculado sobre las variables continuas originales. \\
\textbf{Prueba de permutación (seleccionada)} & No exige ningún supuesto distribucional: solo requiere que, bajo la hipótesis nula de no asociación, cualquier reasignación de los valores de contratación entre municipios sea igual de probable. Construye la distribución nula \textbf{directamente desde los datos observados}, respetando su asimetría real y sus empates, y produce un valor $p$ válido incluso en muestras finitas (no depende de aproximaciones asintóticas). & Es el método estándar para pruebas de independencia libres de distribución cuando los supuestos paramétricos son dudosos (Good, 2005); además permite usar el mismo estadístico de correspondencia definido en la Fase 7.4, en lugar de forzar el uso de una prueba genérica. \\
\bottomrule
\end{longtable}
\end{center}

\textbf{Pasos de ejecución.}
\begin{enumerate}
\item Calcular el estadístico de asociación o correspondencia observado entre
IVS y el indicador de contratación de interés (correlación de Spearman o el
índice de correspondencia definido en la Fase 7.4).
\item Generar $P \geq 5000$ permutaciones de la asignación municipio--valor de
contratación, rompiendo el vínculo real entre vulnerabilidad y contratación
pero preservando la distribución marginal de cada variable.
\item Recalcular el estadístico en cada permutación para construir la
distribución nula empírica.
\item Calcular el valor $p$ empírico como la proporción de estadísticos
permutados iguales o más extremos que el observado.
\end{enumerate}

\textbf{Entregables.} Distribución nula empírica graficada junto al
estadístico observado; valor $p$ empírico reportado junto al valor $p$
paramétrico de la prueba clásica correspondiente, señalando concordancia o
discrepancia entre ambos.

\textbf{Ubicación en el flujo del proyecto.} Se ejecuta como complemento
metodológico dentro de la Fase 7.1 (Análisis de correlación bivariada) y la
Fase 7.4 (Medidas de correspondencia).

\subsection{Línea de trabajo futuro: diseño muestral para validación externa}

Como línea de trabajo futuro --- coherente con la agenda de investigación de
la Sección 11 --- se deja documentado, sin implicar su ejecución dentro del
alcance temporal actual, el diseño conceptual de un muestreo probabilístico
por conglomerados (municipios como conglomerados, estratificados por región y
nivel de vulnerabilidad) para una eventual encuesta de percepción de equidad en
la contratación pública dirigida a funcionarios municipales o ciudadanía. Este
diseño serviría como validación externa cualitativa de los hallazgos
cuantitativos del proyecto, pero su levantamiento excede los recursos y el
tiempo disponibles para el presente trabajo de grado.

\textbf{Justificación: por qué conglomerados estratificados y no muestreo
estratificado simple.} A diferencia de la auditoría de SECOP II (donde la
unidad de muestreo es un contrato y no existe costo de desplazamiento), aquí la
unidad es un municipio que hay que visitar o contactar directamente. Un diseño
estratificado simple podría seleccionar municipios dispersos por todo el país,
lo que encarecería el levantamiento de forma prohibitiva para un trabajo de
grado. Al usar los municipios como conglomerados dentro de cada estrato
región--nivel de vulnerabilidad, el esfuerzo de campo se concentra
geográficamente sin sacrificar la representación de las distintas regiones y
niveles de vulnerabilidad del país. Es, por tanto, el único diseño que
equilibra representatividad estadística con factibilidad logística para este
componente específico, aunque --- como ya se indicó --- su ejecución
efectiva queda fuera del alcance temporal actual.

\subsection{Integración de los procedimientos en el flujo de fases del proyecto}

\begin{center}
\small
\begin{longtable}{p{2.4cm}p{4.6cm}p{6.3cm}}
\toprule
\textbf{Fase del flujo} & \textbf{Procedimiento} & \textbf{Función} \\
\midrule
\endfirsthead
\toprule
\textbf{Fase del flujo} & \textbf{Procedimiento} & \textbf{Función} \\
\midrule
\endhead
2 $\rightarrow$ 3 & Muestreo estratificado de auditoría (Procedimiento 1) & Estimar con IC la tasa de errores de calidad en SECOP II sin auditoría exhaustiva \\
4.6 & \textit{Bootstrap} del IVS (Procedimiento 2) & Cuantificar estabilidad de cargas factoriales y del índice \\
7.1 / 7.4 & Prueba de permutación (Procedimiento 3) & Inferencia no paramétrica complementaria sobre la correspondencia \\
Agenda futura & Diseño muestral para validación externa & Documentar línea de trabajo futura, no ejecutada en este alcance \\
\bottomrule
\end{longtable}
\end{center}

% ================================================================
\section{Limitaciones previstas}

\begin{itemize}
\item El IPM municipal no tiene necesariamente frecuencia anual comparable con
la contratación anual.
\item SECOP II es información administrativa autodeclarada y puede contener
errores, cambios de campos, duplicados o registros incompletos --- riesgo que
el muestreo estratificado de auditoría (Sección 9.3) permite cuantificar, no
eliminar por completo.
\item El valor de un contrato no equivale necesariamente al gasto
efectivamente ejecutado ni al beneficio social generado.
\item La localización de una entidad contratante no implica que el beneficio
económico del contrato permanezca en el mismo municipio.
\item Una asociación estadística no identifica causalidad.
\item Los municipios presentan tamaños poblacionales y estructuras
institucionales muy diferentes.
\item Los cambios metodológicos de SECOP y DANE deberán documentarse antes de
comparar periodos.
\end{itemize}

% ================================================================
\section{Agenda de investigación inmediata}

\begin{enumerate}
\item Auditar completamente las variables y años del histórico SECOP II del
repositorio (formalizado mediante el Procedimiento 1, Sección 9.3).
\item Construir un diccionario de variables de contratación y clasificar cada
campo por función analítica.
\item Identificar las variables municipales DANE compatibles con el IPM y la
dimensión temporal.
\item Definir el mecanismo de integración territorial mediante DIVIPOLA.
\item Realizar EDA y auditoría de calidad antes de construir el índice.
\item Evaluar la factibilidad estadística del PCA y construir versiones
alternativas del índice, validadas mediante \textit{bootstrap} (Procedimiento
2, Sección 9.4).
\item Definir formalmente la medida de correspondencia y compararla con varias
alternativas, incluyendo la prueba de permutación (Procedimiento 3, Sección
9.5).
\item Construir medidas de brecha territorial.
\item Diagnosticar dependencia espacial antes de utilizar modelos espaciales.
\item Definir el modelo temporal/panel según la disponibilidad real de datos.
\item Ampliar la revisión bibliográfica en Scopus, Web of Science, SciELO,
Redalyc y repositorios nacionales.
\item Redactar problema, justificación, objetivo general, objetivos
específicos y metodología definitiva.
\end{enumerate}

% ================================================================
\section{División de trabajo entre integrantes}

\subsection{Principio de división}

La división se organiza alrededor de los dos pipelines de datos independientes
que el proyecto requiere integrar (vulnerabilidad DANE/IPM y contratación
SECOP II), de manera que cada integrante sea responsable completo de una
fuente hasta el punto de integración territorial, evitando duplicar esfuerzo.
Las fases posteriores a la integración (6 a 10) se dividen por tarea
específica dentro de cada fase, dado que dependen de ambas fuentes
simultáneamente.

\begin{center}
\small
\begin{tabular}{ll}
\toprule
\textbf{Integrante} & \textbf{Eje principal} \\
\midrule
Juan Diego Gallo Quintero & Vulnerabilidad socioeconómica (DANE / IPM) \\
Juan David Parada Fonseca & Contratación pública (SECOP II) \\
\bottomrule
\end{tabular}
\end{center}

\subsection{Responsable: Juan Diego Gallo Quintero --- Eje de Vulnerabilidad}

\begin{recuadro}{Fase 2 (parte DANE/IPM) --- Revisión de datos y fuentes}
\begin{enumerate}
\item Auditar las variables disponibles del IPM municipal 2018 (dimensiones,
privaciones, factores de expansión).
\item Verificar cobertura municipal completa e identificar municipios con
datos faltantes o no certificados.
\item Construir el diccionario de variables IPM con definición, escala y
fuente de cada indicador.
\item Evaluar qué otras variables municipales (TerriData/DNP u otras) son
compatibles en periodo y calidad para complementar el IPM.
\end{enumerate}
\textit{Entregable:} diccionario de variables DANE/IPM, matriz de
disponibilidad municipal, decisión documentada sobre variables adicionales
candidatas.
\end{recuadro}

\begin{recuadro}{Fase 4 completa --- Construcción del Índice de Vulnerabilidad Socioeconómica}
\begin{enumerate}
\item Seleccionar y estandarizar los indicadores base (IPM + complementarios
aprobados).
\item Ejecutar pruebas de adecuación muestral (KMO y esfericidad de Bartlett)
antes de correr el PCA.
\item Ejecutar el Análisis de Componentes Principales y decidir el número de
componentes a retener, justificando el criterio (varianza explicada, criterio
de Kaiser, gráfico de sedimentación).
\item Verificar signos de las cargas e invertir variables cuyo aumento
represente menor vulnerabilidad.
\item Construir el IVS final combinando los componentes retenidos.
\item Ejecutar la validación mediante \textit{bootstrap} (Procedimiento 2,
Sección 9.4) sobre las cargas factoriales y el índice resultante.
\item Redactar el reporte de validación y sensibilidad del índice.
\end{enumerate}
\textit{Entregable:} IVS calculado por municipio, cargas factoriales con
intervalos de confianza \textit{bootstrap}, reporte de validación.
\end{recuadro}

\begin{recuadro}{Transversal --- Redacción teórica y bibliográfica}
\begin{enumerate}
\item Mantener y ampliar el estado del arte y el marco teórico con la revisión
bibliográfica pendiente en Scopus, Web of Science, SciELO y Redalyc.
\item Redactar planteamiento del problema, justificación, objetivos y
metodología definitiva del artículo final.
\item Consolidar la matriz bibliográfica y las referencias en formato APA 7.
\end{enumerate}
\textit{Entregable:} documento de estado del arte y marco teórico en su
versión final, integrado con los resultados empíricos del proyecto.
\end{recuadro}

\subsection{Responsable: Juan David Parada Fonseca --- Eje de Contratación}

\begin{recuadro}{Fase 2 (parte SECOP II) --- Revisión de datos y fuentes}
\begin{enumerate}
\item Auditar el notebook \texttt{Historico\_de\_Datos.ipynb} y verificar de
forma directa el número de registros, variables y cobertura temporal
reportados.
\item Ejecutar el muestreo estratificado de auditoría de calidad
(Procedimiento 1, Sección 9.3): definir estratos, calcular tamaño de muestra,
extraer la muestra y codificar las dimensiones de calidad.
\item Estimar, con intervalos de confianza, la tasa de duplicados,
inconsistencias territoriales, campos incompletos y valores atípicos.
\item Definir, junto con Juan Diego, el periodo analítico definitivo a partir
de la calidad y cobertura real encontrada.
\end{enumerate}
\textit{Entregable:} ficha técnica del diseño muestral, reporte de calidad de
SECOP II con intervalos de confianza, recomendación de ventana temporal
definitiva.
\end{recuadro}

\begin{recuadro}{Fase 3 completa --- Preparación y limpieza de datos}
\begin{enumerate}
\item Descargar/integrar los datos SECOP II desde los archivos Parquet
consolidados.
\item Estandarizar entidades territoriales mediante códigos DIVIPOLA,
resolviendo homologaciones y municipios no identificables.
\item Aplicar las reglas de limpieza priorizadas según el resultado del
muestreo de auditoría (Procedimiento 1): tratamiento de duplicados, valores
faltantes y atípicos.
\item Agregar los contratos a nivel panel municipio--año.
\item Construir la base final ``SECOP II limpia'' y el panel municipio--año de
contratación pública.
\end{enumerate}
\textit{Entregable:} base SECOP II limpia, panel municipio--año, métricas de
contratación agregadas.
\end{recuadro}

\begin{recuadro}{Fase 5 completa --- Índices e indicadores de contratación pública}
\begin{enumerate}
\item Calcular métricas principales: valor total, número de contratos, valor
per cápita (usando población municipal como control).
\item Construir indicadores por modalidad de contratación.
\item Construir indicadores por sector.
\item Construir indicadores por tipo de contrato.
\item Construir indicadores por duración (días entre inicio y terminación).
\item Construir los índices compuestos de contratación que resuman las
familias anteriores.
\end{enumerate}
\textit{Entregable:} métricas y clasificación de contratación por
municipio-año, índices compuestos, tablas descriptivas documentadas.
\end{recuadro}

\subsection{Trabajo conjunto --- Fases 6 a 10}

Estas fases dependen simultáneamente del IVS (eje Juan Diego) y de las
métricas de contratación (eje Juan David), por lo que no se dividen por fase
completa sino por tarea específica dentro de cada una. Ambos revisan y validan
el resultado del otro antes de avanzar a la siguiente fase.

\begin{center}
\small
\begin{longtable}{p{1.6cm}p{6.3cm}p{4.5cm}}
\toprule
\textbf{Fase} & \textbf{Tarea} & \textbf{Responsable principal} \\
\midrule
\endfirsthead
\toprule
\textbf{Fase} & \textbf{Tarea} & \textbf{Responsable principal} \\
\midrule
\endhead
6 & Estadísticos descriptivos y distribuciones por municipio/año & Ambos (cada uno su eje) \\
6 & Mapas temáticos (IVS y contratación) y curvas de Lorenz/Gini & Juan Diego \\
6 & Comparaciones departamentales de contexto & Juan David \\
7 & Análisis de correlación bivariada + prueba de permutación (Procedimiento 3) & Juan David \\
7 & Modelos de regresión múltiple y de panel & Juan Diego \\
7 & Definición de la medida de correspondencia y de brechas territoriales & Ambos (definición conjunta) \\
8 & Matriz de pesos espaciales y pruebas de autocorrelación (Moran's I, LISA) & Juan Diego \\
8 & Interpretación conjunta de clústeres espaciales & Ambos \\
9 & Reducción de dimensionalidad y clustering de municipios & Juan David \\
9 & Perfilamiento socioeconómico y contractual de los clústeres & Ambos \\
10 & Diseño de dashboard y mapas interactivos & Juan David \\
10 & Redacción del reporte final y recomendaciones de política pública & Juan Diego \\
\bottomrule
\end{longtable}
\end{center}

\subsection{Regla de coordinación}

\begin{itemize}
\item Ningún integrante avanza a una fase conjunta (6 en adelante) sin que los
entregables de su fase individual (2--5) estén validados por el otro
integrante.
\item La integración territorial (unión de los paneles de vulnerabilidad y de
contratación por código DIVIPOLA--año) se ejecuta de forma conjunta, dado que
requiere acuerdo explícito sobre el nivel de agregación y el tratamiento de
casos no homologables.
\item Todo cambio en el periodo analítico definitivo, en la definición de
correspondencia o en el tratamiento de valores atípicos debe documentarse y
acordarse entre ambos antes de ejecutarse, dado su impacto sobre el trabajo del
otro integrante.
\end{itemize}

% ================================================================
\section{Nota metodológica sobre originalidad}

El repositorio revisado es una referencia externa utilizada como punto de
partida para identificar fuentes, estructura de datos y antecedentes. El
presente proyecto se plantea como una investigación independiente. El aporte
original propuesto consiste en integrar, a escala municipal, un índice de
vulnerabilidad socioeconómica basado en IPM/PCA con medidas de contratación
pública, y desarrollar un marco estadístico explícito --- que incluye
muestreo, remuestreo y pruebas de permutación --- para estudiar asociación,
correspondencia y brechas, incorporando dimensión temporal y espacial cuando
los datos lo justifiquen.

% ================================================================
\section{Referencias bibliográficas --- APA 7}

Arachia, G., Assisi, D., Cesi, B., Giuranno, M. G., \& Russo, F. (2024).
Intermunicipal cooperation in public procurement. \textit{Regional Studies,
58}(11), 2055--2073. https://doi.org/10.1080/00343404.2024.2350615

Ardila-Rueda, L. (2004). Gasto público y convergencia regional en Colombia.
\textit{Ensayos Sobre Política Económica, 22}(45), 222--268.

Bonet-Morón, J. (2006). \textit{Desequilibrios regionales en la política de
descentralización en Colombia} (Documento de Trabajo sobre Economía Regional y
Urbana No.\ 77). Banco de la República.

Cochran, W. G. (1977). \textit{Sampling techniques} (3.\textsuperscript{a}
ed.). John Wiley \& Sons.

Colombia Compra Eficiente. (2026). \textit{SECOP II: Sistema Electrónico para
la Contratación Pública}.

DANE. (2018). \textit{Medida de pobreza multidimensional de fuente censal}.
Departamento Administrativo Nacional de Estadística.

Datos Abiertos Colombia. (2026). \textit{SECOP II --- Contratos Electrónicos}.

Departamento Nacional de Planeación. (2024). \textit{TerriData: Sistema de
Estadísticas Territoriales}. DNP.

Díaz Díez, C. A. (2023). Obligatoriedad del Sistema Electrónico de
Contratación Pública (SECOP) para las entidades estatales con régimen
contractual excepcional: un supuesto de retorno del Derecho Administrativo en
Colombia. \textit{Revista Eurolatinoamericana de Derecho Administrativo,
10}(2).

Efron, B., \& Tibshirani, R. J. (1993). \textit{An introduction to the
bootstrap}. Chapman \& Hall/CRC.

Ferry, L., et al. (2023). Public procurement as a policy tool: The territorial
dimension. \textit{Regional Studies, 57}(11), 2087--2101.
https://doi.org/10.1080/00343404.2022.2134850

Figueroa-Gómez, Y., \& Galpin, I. (2024). Identifying public tenders of
interest using classification models: A comparative analysis. \textit{SN
Computer Science, 5}, 87. https://doi.org/10.1007/s42979-023-02443-3

Good, P. I. (2005). \textit{Permutation, parametric, and bootstrap tests of
hypotheses} (3.\textsuperscript{a} ed.). Springer.

Gutiérrez Vanegas, K. F. (2024). \textit{Implementación de banderas rojas en
la contratación pública en Colombia utilizando datos abiertos} [Trabajo de
grado]. Fundación Universitaria Los Libertadores.

Lohr, S. L. (2019). \textit{Sampling: Design and analysis}
(3.\textsuperscript{a} ed.). CRC Press.

Monsalvo Herrera, A. E., \& Jiménez, W. G. (2025). Pobreza multidimensional:
un análisis espacial en los municipios de Cundinamarca, Colombia.
\textit{Novum Jus, 19}(2), 213--251.
https://doi.org/10.14718/NovumJus.2025.19.2.8

Norbäck, M., \& Zapata Campos, M. J. (2022). The market made us do it: Public
procurement and collaborative labour market inclusion governance from below.
\textit{Social Policy \& Administration, 56}(4), 632--647.
https://doi.org/10.1111/spol.12787

Open Contracting Partnership. (2023). \textit{Open contracting reforms to
boost civic engagement and participation: The experience of nine regions in
Colombia}.

Rodriguez-Plesa, E., Dimand, A.-M., \& Alkadry, M. G. (2022). Community social
capital, political values, or organizational capacity? Indicators of
engagement in sustainable public procurement at the local level.
\textit{Journal of Cleaner Production, 338}, 130556.
https://doi.org/10.1016/j.jclepro.2022.130556

Salazar, A., Pérez, J. F., \& Gallego, J. (2024). VigIA: Prioritizing public
procurement oversight with machine learning models and risk indices.
\textit{Data \& Policy, 6}, e75. https://doi.org/10.1017/dap.2024.83

Sarter, E. K. (2024). Public procurement and labour market inequality:
Conceptualising a multi-faceted relationship. \textit{Journal of Industrial
Relations, 66}(1), 130--152. https://doi.org/10.1177/00221856231221417

Särndal, C.-E., Swensson, B., \& Wretman, J. (1992). \textit{Model assisted
survey sampling}. Springer-Verlag.

\end{document}
